% Paper 1 manuscript skeleton, gradvar-phoenix programme (SSIT Tumakuru).
% Status: SKELETON. No hardware minute has been spent. Every quantity marked [RESULT], [TBD] or \placeholder{}
% is a placeholder to be filled from the logged data after the post-run review; nothing here is a measured result.
% Pre-registration: preregistration_q1.html v0.11.1 (19-20 Sep 2026, Deviations 14-42). Analysis follows Sections 3 and 3b exactly as registered.
% Build: latexmk -pdf main.tex (see ../Makefile). Class: REVTeX 4.2, PRX Quantum / QST style.
\documentclass[aps,prxquantum,reprint,superscriptaddress,nofootinbib,longbibliography,amsmath,amssymb]{revtex4-2}
\usepackage{graphicx}
\usepackage{booktabs}
\usepackage{longtable}
\usepackage{bbm}
\usepackage{xcolor}
\usepackage{hyperref}
\hypersetup{colorlinks=true,linkcolor=blue!60!black,citecolor=blue!60!black,urlcolor=blue!60!black}

% --- placeholder macros: every result field goes through one of these so a grep for RESULT finds them all
\newcommand{\placeholder}[1]{\textcolor{red!70!black}{[RESULT: #1]}}
\newcommand{\tbd}[1]{\textcolor{orange!80!black}{[TBD: #1]}}
\newcommand{\figstub}[2]{\fbox{\parbox[c][0.28\textheight][c]{0.95\linewidth}{\centering\textcolor{gray}{Figure stub: \texttt{#1}\\[1ex]#2}}}}
\newcommand{\Var}{\mathrm{Var}}
\newcommand{\Cmix}{C_{\mathrm{mix}}}
\newcommand{\Np}{\mathcal{N}_p}
\newcommand{\eN}{\varepsilon_N}

\begin{document}

\title{Pre-registered measurement of parameter-shift gradient variance on a 120-qubit superconducting processor with a controlled non-unital reset dial}

\author{\tbd{Author 1}}
\affiliation{Department of Computer Science and Engineering, SSIT Tumakuru, India}
\author{\tbd{Author 2}}
\affiliation{Department of Computer Science and Engineering, SSIT Tumakuru, India}
\author{\tbd{Physics co-author, if named at submission}}
\affiliation{\tbd{affiliation}}

\date{\today}

\begin{abstract}
Barren-plateau theory predicts how the variance of a parameter-shift gradient of a local observable scales with qubit count $n$ and depth $L$ for a hardware-efficient ansatz, under noiseless, unital and non-unital noise. We report a pre-registered hardware test of these predictions on \texttt{ibm\_phoenix} (IBM Nighthawk r2, 120 qubits on a $12\times10$ square lattice): a Ry/CZ ansatz on rectangular patches of $n = $ \placeholder{actual $n$ ladder, e.g.\ 20/39/53/70/87} qubits at $L \in \{1, 2, 4, 8, 12\}$, with $M = 200$ uniform parameter draws per point, $4096$ shots per shifted circuit and a two-term shift rule, every hypothesis, threshold and kill rule fixed before the first allocation minute. Predictions for every point were drawn before data by exact light-cone simulation and, at $L \geq 8$, by second-moment Pauli propagation with a rigorous truncation lower bound. A reset applied to each qubit with probability $p$ after every layer realises the non-unital mixture channel $\Np(\rho) = (1-p)\rho + p|0\rangle\langle 0|$ and turns the device's native reset into a dial on the non-unital strength, tested at $p \in \{0.25, 0.5\}$ against delay-matched and unital-dephasing controls. We find that the depth dependence at fixed $n$ is \placeholder{exponential/power-law fit verdict and exponent with 95\% bootstrap interval}, the $n$ dependence at fixed $L$ is \placeholder{verdict per $L$}, the noiseless-normalised layer-index ratio is \placeholder{value and interval against the unital and non-unital simulations}, and under the reset dial the last-layer gradient variance and the cost variance \placeholder{follow / do not follow} the pre-drawn second-moment curves; the headline statistic, the measured $k = L$ variance over the rigorous ansatz-specific floor $\tfrac12 c_i^2 g_i^2 p^2$ (Deviation 33), is \placeholder{value with interval} against the pre-drawn 1.88 ($p = 0.25$) and 1.30 ($p = 0.5$). The arm is a variance-level benchmark of a first-principles noise model, not a test of the 2-design theorems and not a demonstration of simulability. All raw counts, calibration snapshots and the pre-registration with its Deviations table are public.
\end{abstract}

\maketitle

\section{Introduction}\label{sec:intro}

\paragraph{Barren plateaus.} For random parameterised circuits that form approximate 2-designs, the variance of a cost gradient over the parameter distribution is exponentially small in the qubit count~\cite{McClean2018}; for shallow circuits with local observables the decay with $n$ is polynomial up to a depth that grows logarithmically in $n$ in one dimension~\cite{Cerezo2021}. For deep hardware-efficient circuits the variance is set by the dimension of the circuit's dynamical Lie algebra once the circuit forms a 2-design over that algebra~\cite{Ragone2024,Fontana2024,Larocca2025}; for shallower circuits the Lie-algebraic formula does not hold and the variance can only be bounded~\cite{Larocca2025}. Local Pauli noise induces a plateau of its own: Wang et al.~\cite{Wang2021} bound the derivative itself for all parameters by $\sqrt{8\ln 2}\,N_O\|\omega\|_\infty\|\eta\|_1\,\sqrt{n}\,q^{cL+1}$ with $c = 1/(2\ln 2)$, and Mele et al.~\cite{Mele2026} (Proposition 10) give $\Var[\partial C] \leq e^{-\Omega(|P|+L)}$ for any unital non-unitary noise with no $n$ prefactor. Under any local non-unital noise acting after every layer, on average over circuits whose two-qubit gates form local 2-designs, the same authors prove that local-observable costs keep a $\Theta(1)$ variance in $n$ with lower bound $(\|t\|_2^2/3)^{|P|}$ (Corollary 6), where $t$ is the Bloch-translation vector of the per-layer channel, and that only parameters in the last $\Theta(\log n)$ layers have non-vanishing gradient variance, $\Var[\partial_k C] = \exp(-\Theta(L-k))$ (Theorem 8)~\cite{Mele2026}; Singkanipa and Lidar~\cite{Singkanipa2025} treat the same non-unital regime and give a noisy parameter-shift rule for Hilbert--Schmidt-contractive channels. The absence of a plateau in this regime is not by itself good news: it arises because the circuit becomes effectively shallow~\cite{Mele2026}, and average-case classical simulation of such noisy circuits by Pauli propagation has a proven guarantee~\cite{Angrisani2026}. The review of Larocca et al.~\cite{Larocca2025} lists the relationship among noise, measurement and concentration as open, and its only methodological guidance on measuring plateaus (the ``Methods for analysing variance'' paragraph and Box~2) gives no sample counts or reporting standards.

\paragraph{The gap.} \tbd{Two sentences, to be checked against the final literature sweep before submission:} no hardware barren-plateau study has reported the variance of a parameter-shift gradient over random parameters with confidence intervals on more than a handful of physical qubits (the only published hardware parameter-shift error measurement is the five-qubit mean-squared-error test of Mari et al.~\cite{Mari2021}, whose hardware floor sat an order of magnitude above the analytic shot floor), and none has pre-registered its hypotheses, thresholds and kill rules before data. The non-unital predictions in particular have not been tested where they are resolvable: at the native $T_1$ of a superconducting device the per-layer translation norm is $\|t\|_2 \approx 2\times10^{-3}$ and the proven floor for a weight-2 observable is of order $10^{-12}$, below any shot floor.

\paragraph{This work.} We measure $\Var_\theta[\partial C]$ for one parameter-shift gradient of $Z_iZ_j$ under a Ry/CZ ansatz on \texttt{ibm\_phoenix} across a ladder in $n$ and $L$, with pre-drawn predictions from three noise models, and add a controlled non-unital arm in which the device's native 400\,ns reset, applied with probability $p$ per qubit per layer and pooled over random masks, realises $\Np$ at $p \in \{0.25, 0.5\}$. Section~\ref{sec:methods} gives the device, ansatz, estimator, pre-registration and gates; Section~\ref{sec:prereg} restates the pre-registered hypotheses and refutation conditions verbatim; Section~\ref{sec:results} reports every registered statistic in tables keyed to the job-list point identifiers; Section~\ref{sec:discussion} interprets each outcome branch with text written before data, as the pre-registration allows.

\section{Methods}\label{sec:methods}

\subsection{Device}\label{sec:device}
\texttt{ibm\_phoenix} is an IBM Nighthawk r2 processor with 120 qubits on a $12\times10$ square lattice (218 couplers, CZ native). On the planning snapshot of 19 September 2026 the median CZ error was $1.9\times10^{-3}$, the median readout error $8.2\times10^{-3}$, $T_1$ about 178\,$\mu$s, the CZ duration 68\,ns (208\,ns maximum), single-qubit gates 40\,ns, measurement 1.94\,$\mu$s and the native reset 400\,ns on 119 qubits (2140\,ns on qubit 79). The device default repetition delay is 1.0\,$\mu$s (read from \texttt{backend.configuration()} on 20 September 2026, Deviation 23). Run-day values: \placeholder{table of run-day medians per hardware day from \texttt{data/calibrations/}}.

\subsection{Ansatz, observable and patches}\label{sec:ansatz}
One layer is $R_y(\theta)$ on every qubit of the patch followed by CZ on every lattice edge of the patch in four sub-layers (horizontal even, horizontal odd, vertical even, vertical odd); $L$ layers give $nL$ parameters, drawn uniformly from $[0, 2\pi)$. Uniform draws characterise the landscape and are not a training recommendation~\cite{Larocca2025}. The observable is $Z_iZ_j$ on one interior edge of the patch, fixed from the snapshot for the whole campaign. Patches are rectangles $4\times5$, $4\times10$, $6\times10$, $8\times10$, $10\times10$ placed inside the clean component by the exclusion rules of Section~\ref{sec:exclusion}; couplers with CZ error $\geq 5\times10^{-3}$ are treated as broken lattice edges carrying no CZ (Deviation 26). The observable edge must be an intact coupler with no broken coupler within graph distance 2 where the patch allows it. On the 19 September snapshot the ladder is $n = 20/39/53/70/87$ with observable edges 93\_103, 93\_103, 84\_85, 84\_85, 75\_85 (Deviations 18 and 36; the $4\times10$ edge moved from 94\_95 because the broken coupler (95, 96) sat on an observable qubit, and the $n = 39$ Gate 1 and Gate 1b rows are to be recomputed on the new edge before the Gate 1 decision); the run-day placements are \placeholder{table: patch, origin, holes, broken couplers, edge, $n$ per hardware day, from \texttt{layout\_check}}.

\paragraph{Light cone versus patch.} The Heisenberg support of $Z_iZ_j$ grows by graph distance one per layer, so the causal cone at depth $L$ is the radius-$(L-1)$ neighbourhood of the edge, about $2L^2$ qubits: contained in every patch at $L = 1, 2$; contained at $L = 4$ only by patches of at least $7\times7$ (the $8\times10$ and $10\times10$ rungs); truncated by every patch at $L = 8, 12$. The noiseless variance depends on $n$ only through the cone graph, the qubits and intact couplers within distance $L - 1$ of the edge after the patch boundary, the excluded qubits and the broken couplers are removed~\cite{Uvarov2021}. It is exactly $n$-independent at $L = 1$ on every patch; at $L \geq 2$ it is $n$-independent only between rungs whose cone graphs are isomorphic (on the 19 September placements the $4\times5$ and $4\times10$ $L = 2$ cones coincide, 16 qubits and 24 CZ, as do the $6\times10$ and $8\times10$ cones, 16 qubits and 22 CZ; the $10\times10$ cone matches no other rung, and from $L = 3$ every cone contains holes or broken couplers; Deviation 36). Where cone graphs differ, a same-geometry pre-drawn curve tests the noise model on that graph, not $n$ dependence, and is reported as such; the cone graph of every hardware point is logged and simulated as placed. Where the cone is truncated the $n$ dependence is a boundary effect, and the 4-row patches sit systematically above the larger ones at $L \geq 3$ for geometric reasons.

\subsection{Gradient estimator}\label{sec:estimator}
The two-term parameter-shift rule with shifts $\pm\pi/2$ is exact for any gate with an involutory generator, so for $R_y$~\cite{Mari2021}; it is exact for the noisy expectation value provided the noise channels do not depend on $\theta$, which holds for virtual $R_z$~\cite{Banchi2026}. The differentiated parameter is $\theta_{(k,q)}$ with $q$ one of the observable's qubits (the cone centre) and $k$ the layer index; the main grid uses $k = 1$. Per draw we run three circuits ($\theta \pm \pi/2$, $\theta$) at 4096 shots (16384 at the $n = 100$, $L = 8$ headline points); the unshifted circuit gives $C(\theta)$, and at 5 draws per point a fourth circuit at $\theta + \pi$ tests the rotation-survival identity $f(\theta+\pi) = f(\theta+\pi/2) + f(\theta-\pi/2) - f(\theta)$~\cite{Mari2021}, which holds only if $\theta$ enters a single rotation after transpilation (logged as \texttt{param\_expr}). The estimator is unscaled and unbiased, so $\mathbb{E}[\hat g^2] = \Var_\theta[g] + \mathbb{E}[\text{shot variance}]$, and rescaling would not improve distinguishability~\cite{AghaeiSaem2026}. The shot-noise variance of one gradient estimate is $[\sigma_0^2(\theta+\pi/2)+\sigma_0^2(\theta-\pi/2)]/(4N) \approx \sigma_0^2/(2N)$ with $\sigma_0^2 \leq 1$ the single-shot variance of $Z_iZ_j$ and $N$ the shots per shifted circuit~\cite{Mari2021}; we use the exact two-term form with the measured spreads. Floors: $1/(2\cdot4096) = 1.22\times10^{-4}$ and $1/(2\cdot16384) = 3.05\times10^{-5}$.

\paragraph{Deviation 25 estimator (simulation regression).} The pre-registered regression check, $\Var = 2^{-n}$ for the one-layer chain with observable $Z^{\otimes n}$ inside a percentile-bootstrap interval, cannot be met at $n \geq 13$ at any affordable $M$ because the chain gradient $g = -\sin\theta_0\prod_{i>0}\cos\theta_i$ has kurtosis $(3/2)^n$. The amended check (Deviation 25) is (a-i) a per-draw closed-form identity $|g_{\mathrm{sim}} - g_{\mathrm{closed}}| < 10^{-10}$ at $n = 4..20$, (a-ii) bootstrap coverage at $n \leq 12$ with $M = 1000$, (a-iii) the sample variance inside the central 95\% of its exact null sampling distribution at $n = 13..20$ with $M = 300$, with a seed-2027 replicate rule. Status before data: (a-i) max $1.8\times10^{-14}$, (a-ii) 9/9, (a-iii) 7/8 at seed 2026, the $n = 20$ point at the 2.4th percentile resolved by the seed-2027 replicate (0.212 inside [0.043, 5.764]) (Gate 1 report, Section 2). The Arrasmith et al.\ closed forms $(1/8)(3/8)^{n-1}$ and $1/(8n^2)$~\cite{Arrasmith2022} are the second regression reference.

\subsection{Pre-registration, gates and review rule}\label{sec:gates}
The pre-registration (v0.1 signed 19 September 2026; v0.11.1 the version at booking) fixes hypotheses H1--H7, the grid, the analysis plan and the kill rules before any allocation minute. Four gates are decided in order before the corresponding minutes are booked: Gate 1 (simulation: criteria (a)--(f) on the pre-drawn predictions), Gate 1b (pre-booking resolvability of the dial arm's ladder by Pauli propagation), the smoke test and Gate 2 (hardware: variance above the null-control floor at $n = 20$, budget, readout and reset error within $1.5\times$ of planning). An independent reviewer signs off the exact job list before every submission and audits the logged rows against the pre-registration before any fit is run. Every change after signature is a numbered Deviation with date, reason and approver (Appendix~\ref{app:deviations}); Deviations 14--42 were all recorded before any Paper 1 datum (Deviation 32 fixes, before the smoke test's data existed, that they feed Gate 2 and the kill rules only and never a hypothesis); Gate 1b clause (b) was amended three times before data (Deviations 28--30), frozen, and then clarified without change of direction or power (Deviation 35: resolvability of a rung's reference is judged on the measured, floor-subtracted value, never on the prediction; the clause passes only if at least two rungs are counted and all counted rungs pass, and is inconclusive otherwise; the $n = 53$, $L = 8$, $p = 0$ reference runs at 16384 shots so that all three rungs can count), a history we state plainly. Deviations 33--42 adopt the delegated theorist referee pass as modified by its second review and the analysis-pipeline clarifications; the ten Deviation 42 readings are the ones the analysis code implements. The \texttt{ibm\_phoenix} smoke test ran on 20 September 2026 (Deviation 32; job list 02, about 0.7 Flex minutes at the 1\,$\mu$s default from the dry-run line). Gate outcomes: \placeholder{Gate 1 verdict and date; Gate 1b; Gate 2 with the smoke-test numbers}.

\subsection{Budget model and execution}\label{sec:budget}
Locked QPU time per job is modelled (Deviation 24) as $T = 2\,\mathrm{s} + (\tau_{\mathrm{rep}} + L\times0.71\,\mu\mathrm{s} + t_{\mathrm{meas}} + 10\,\mu\mathrm{s}\,[+\text{dial durations}])\,N_{\mathrm{exec}}\,(3\text{ at resilience 2}) + [\text{resilience}\geq1]\,32\,\text{shots}\,n_{\mathrm{bases}}\,(\tau_{\mathrm{rep}} + t_{\mathrm{meas}} + 10\,\mu\mathrm{s})$, fitted on an \texttt{ibm\_marrakesh} pipeline check (21.1\,s predicted against 22\,s charged). At the booked $\tau_{\mathrm{rep}} = 1\,\mu$s the Paper 1 grid (about $2.85\times10^{8}$ executions, 188 jobs) is about 90\,min against a 200-minute line; the dial core is 57.8\,min against 65; the programme ledger is $200 + 65 + 45 + 50 = 360$ Flex minutes. Measured: \placeholder{locked time per point and per job from \texttt{qpu\_seconds}, against the model}. Execution used EstimatorV2 at resilience 0 and 1 on the full grid and 2 (linear ZNE from gains 1 and 3, $c = (3/2, -1/2)$, $\|c\|_1 = 2$) on the reduced grid; fractional gates disabled; the granted $\tau_{\mathrm{rep}}$ and \texttt{dynamic\_reprate\_enabled} logged per job.

\subsection{Exclusion and layout rules}\label{sec:exclusion}
Qubits are excluded if dead (qubit 17), in the high-error CZ cluster (55, 61, 62, 63, 72, 73 on the snapshot; the 19 September instance of the coupler cut), readout error $> 3\times10^{-2}$ on the day, $|ZZ| \geq 1$\,MHz to an excluded or dead qubit, or initialisation error $\geq 5\times10^{-4}$ (the last two from the raw backend properties, Deviation 22). Before every submission the runner re-checks the placed patch against the live properties (readout, CZ $< 5\times10^{-3}$ per coupler, initialisation, ZZ) and logs the verdict, any one-column patch shift and any override (never on the observable edge or its $L = 2$ cone) in \texttt{layout\_check} (Deviation 26). Run-day verdicts: \placeholder{count of pass/shift/override per job}.

\subsection{Pre-drawn predictions}\label{sec:predictions}
Classical baselines beside every hardware point: exact statevector or density-matrix simulation on the light cone up to 24 cone qubits with 32 noise trajectories; second-moment Pauli propagation at $L \geq 8$ and for the large-cone $L = 4$ points (Deviation 15); noisy models with the snapshot's per-edge CZ error, per-qubit 1Q error, $T_1$, $T_2$, readout confusion and, in both the main-grid and the dial models, static $ZZ$ on every lattice pair with a reported coupling over the whole scheduled layer outside calibrated CZs (interaction $e^{-i\zeta\tau Z_aZ_b/4}$ with $\zeta = 2\pi\times$ the reported coupling, median 27\,kHz, propagated by the incoherent rule that moves weight $\sin^2(\zeta\tau/2)$ per pair per layer, $8.8\times10^{-4}$ to $3.6\times10^{-3}$ at the median; Deviation 34). The unital model sets the Bloch translation $t = 0$; the non-unital model carries thermal relaxation per CZ sub-layer. Pauli propagation runs two engines on one op program: a deterministic truncation by coefficient whose result is a rigorous lower bound $V_{\mathrm{trunc}} \leq V$, and an unbiased Pauli-path Monte Carlo with a standard error; the prediction is $V_{\mathrm{MC}} \pm 2\sigma$ and the Deviation 15 error is $\max(2\sigma, V_{\mathrm{MC}} - V_{\mathrm{trunc}})$ (Appendix~\ref{app:pauliprop}). Validation: 46 of 48 rows inside the 95\% interval of exact $M = 200$ estimates, the two misses within 5\% of the interval width. Pre-drawn values are quoted from \texttt{data/predictions/} at the commit recorded in Table~\ref{tab:provenance}.

\subsection{Non-unital arm: reset dial}\label{sec:dial}
A reset applied with probability $p$ to each qubit after each layer, averaged over the reset lottery, is the mixture channel $\Np(\rho) = (1-p)\rho + p|0\rangle\langle0|$: Bloch form $r \mapsto Dr + t$ with $D = (1-p)\mathbbm{1}$, $t = (0, 0, p)$, contraction parameter $c = (\|D\|_F^2 + \|t\|^2)/3 = (1-p)^2 + p^2/3$ (0.583 at $p = 0.25$, 0.333 at $p = 0.5$). A fresh mask is drawn per layer and per circuit, the observable's qubits in the lottery like every other qubit; the two shift circuits of a draw share their mask sequence or draw independent ones, fixed per point before submission and logged (Deviation 38); $K = 256$ masks $\times$ 16 shots per (draw, shift) (Deviation 27); estimator $\Cmix(\theta) = K^{-1}\sum_i C_{\mathrm{mask}\,i}(\theta)$, whose expectation is the cost of the circuit with $\Np$ after every layer. Floors: shot floor $1/(2\cdot4096)$ on the gradient and $1/4096$ on $\Var[\Cmix]$; pattern-noise floor $\Var_{\mathrm{mask}}[C]/K$ on $\Var[\Cmix]$ (exact: $\Var_\theta[\Cmix] = \Var_\theta[C_{\Phi_{\mathrm{mix}}}] + \mathbb{E}_\theta\Var_{\mathrm{mask}}[C]/K$), and on the gradient a floor estimated directly from the $K$ per-mask gradient differences (their variance over masks minus their shot variance, divided by $K$), of which $\Var_{\mathrm{mask}}[C]/(2K)$ is the independent-mask upper bound (shared masks cancel most of the last-layer lottery in the difference; Deviation 38); uncertainty propagated by bootstrap over masks within draws. Refutation threshold and headline: the rigorous ansatz-specific floor of Deviation 33, proven for this ansatz and channel from uniform angles alone, $\Var[\partial_{k=L}C] \geq \tfrac12 c_i^2 g_i^2 p^2$ and $\Var[C] \geq \tfrac12 p^2(c_i^2 g_i^2 + c_j^2 g_j^2)$, with $c_i = a_i(1-p)(a_j p + b_j)$ from the readout confusion and $g_i$ the last-layer depolarising factor, $L$- and $n$-independent and pre-drawn per patch and per estimator from the run-day snapshot (on the 19 September snapshot at $n = 53$: $1.06\times10^{-3}$ / $7.35\times10^{-3}$ for $k = L$ and $2.20\times10^{-3}$ / $1.50\times10^{-2}$ for $C$ at $p = 0.25$ / $0.5$, which the Gate 1b predictions clear by 1.88 / 1.30 and 1.62 / 1.22). The measured $k = L$ variance divided by this floor is the arm's headline statistic. The Corollary 6 bound of Mele et al.~\cite{Mele2026} (corrected by Deviation 21), $(\|t\|_2^2/3)^{|P|} = p^4/9$, carries the constants of a 2-design ensemble, which Ry/CZ is not~\cite{Uvarov2021}, and is kept on every panel as a historical reference line only. Matched controls at $n = 60$, $L = 8$: delay-matched $p = 0$ (every reset replaced by \texttt{delay(400\,ns)}) and a unital dephasing dial ($Z$ with probability $p/2$ per qubit per layer, $t = 0$). Truncation arm at $p = 0.5$: the full circuit and the circuit with its first $L - \ell$ layers deleted, $\ell \in \{2, 4\}$, on shared $\theta$ and masks. Kill rules (a)--(d) of Section 3b apply (rule (b): 7.0\,min of QPU-locked time per dial gradient point including job overhead, Deviation 41). Characterisation: \placeholder{reset error per dial qubit, X-basis residual, transpiled measure count, rotation-survival pass count}. The same channel lies inside the HS-contractive class of Singkanipa and Lidar~\cite{Singkanipa2025} and inside the noise class of Angrisani et al.~\cite{Angrisani2026}. The arm is a variance-level benchmark of a first-principles noise model: the floor is the part a reset creates trivially in the last layer and the excess over it is the model-dependent, deeper-path content. Nothing here is a test of Theorem 8 of Mele et al.\ (arXiv Theorem 10), which assumes 2-design layers, nor a demonstration of simulability (Deviation 33).

\section{Pre-registered analysis}\label{sec:prereg}
The text of this section restates the hypotheses and refutation conditions of pre-registration v0.11.1 Sections 1 and 3b; wording changes from the signed v0.1 are Deviations 14, 16, 20, 21, 33, 35--38, 40 and 42 (Appendix~\ref{app:deviations}); the $L = 12$ main-grid rows and the dial $k = 1$ rows at $L = 12$ are exploratory and enter no confirmatory statistic (Deviation 37).

\subsection{Estimands and fits}
Variance of the gradient across the $M$ draws with a 95\% bootstrap interval (10\,000 resamples); the shot-noise contribution subtracted and reported separately. Fits: exponential in $L$ at fixed $n$; exponential in $n$ at fixed $L$; exponential in $L - k$ for the layer-index sweep; a power law fitted alongside each exponential and the better model reported by AIC. The sample mean of $\partial C$ over draws is a null test (exactly zero for uniform draws under $\theta$-independent noise~\cite{Arrasmith2022}). $\Var[C]$ and $\Var[C(\theta_+) - C(\theta_-)]$ are secondary estimands beside $\Var[\partial C]$; the equivalence between plateaus and cost concentration holds up to a factor $m^2L^2$~\cite{Arrasmith2022}, so this classifies exponential against polynomial without giving the exponent. The resolvability ratio $\eN = \Var_{\mathrm{shot}}/(N\Var_\theta)$~\cite{AghaeiSaem2026} is required below 1 for any point to be claimed; ``at any main-grid point'' in the fits and in H4 excludes unclaimable points ($\eN \geq 1$), which are listed with their $\eN$ (Deviation 42 (x)). The smallest signal to be claimed must lie above the measured null-control floor plus 3 bootstrap $\sigma$ of that floor (Deviation 37); rows below that bar are upper bounds. The distribution of $\partial C$ (max $|\partial C|$, kurtosis) is reported per point. All plots log scale with the null-control floor and the analytic shot floor on every panel.

\subsection{H1: noiseless structure}
For $Z_iZ_j$ on the square lattice the noiseless gradient variance is independent of $n$ at fixed $L$ while the patch contains the cone, so exactly $n$-independent at $L = 1$ on every rung and at $L \geq 2$ between rungs with isomorphic cone graphs (Deviation 36); the Uvarov--Biamonte lower bound $3^{-|C|}$ holds for 2-design blocks, which Ry/CZ blocks are not~\cite{Uvarov2021}; on a 2D lattice it is polynomial in $n$ only for $L$ of order $\sqrt{\log n}$~\cite{Cerezo2021}. Once the cone covers the patch and the circuit approaches a 2-design the variance decays exponentially in $n$. \emph{Refuted if} the noiseless simulation at $L = 4$ on patches of at least $7\times7$ shows a variance changing by more than a factor of 2 between $n = 80$ and $n = 100$ outside the bootstrap interval, or if the measured $L = 12$ curve does not fit an exponential in $n$ better than a power law, the simulated curve reported alongside; since $L = 12$ rows are exploratory (Deviation 37) this part of H1 is reported, not decided (Deviation 42 (i)). The $L = 1, 2$ comparisons are a pipeline check only.

\subsection{H2: noise, Pauli type}
Under local Pauli gate noise the measured variance decays exponentially with $L$ at fixed $n$ for every $k$. The Wang et al.\ bound~\cite{Wang2021} is about 23 at $n = 100$, $q = 1 - 2\times10^{-3}$ for every $L \leq 12$ and is uninformative here; the tighter analytic reference is Proposition 10 of Mele et al.~\cite{Mele2026}; the quantitative curve is our unital simulation ($t = 0$) with the global-depolarising sanity model $\Var \propto p^{2L}$ drawn beside it~\cite{KimOz2022}. \emph{Refuted if} any evaluable $(n, k)$ series fails to fall with $L$ once the shot floor is subtracted, or any evaluable $L = 8$ / $L = 12$ $k = L$ pair sits at the same level within its intervals; every evaluated series and pair is listed with its verdict (Deviation 42 (ii)).

\subsection{H3: noise, non-unital (consistency check, Deviation 16)}
Under local non-unital noise after every layer, on average over local 2-design circuits, $\Var[\partial_k C] = \exp(-\Theta(L-k))$ (Theorem 8) and the local cost variance is $\Theta(1)$ with floor $(\|t\|_2^2/3)^{|P|}$ (Corollary 6)~\cite{Mele2026}; on \texttt{ibm\_phoenix} $\|t\|_2 \approx 1.7\times10^{-3}$ and the floor is about $10^{-12}$, below any shot floor. Statistic (Deviation 14): $R_m = [\Var_m(k{=}L)/\Var_m(k{=}1)]/[\text{same, noiseless}]$ on shared draws, $D = R_{\mathrm{nonunital}} - R_{\mathrm{unital}}$. The refutation test is the Deviation 14 clause (the paired-bootstrap interval of $R_{\mathrm{hardware}}/R_{\mathrm{unital}}$ excludes 1 in the predicted direction); Deviation 16's consistency reading ($D = 0$ within the interval) is reported alongside (Deviation 42 (iii)). The ratio is formed only where the $k = 1$ denominator lies above the shot floor; at native noise its predicted value is 0.98--1.05 on converged rows with a $2\sigma$ of 0.04--0.42, so it has no power and is a consistency check (Deviation 40).

\subsection{H4: mitigation}
Readout mitigation (level 1) rescales landscape and shot variance alike, about $(1 - 2p_{\mathrm{ro}})^{-4} \approx 1.07$, and leaves their ratio unchanged~\cite{AghaeiSaem2026}. ZNE (level 2) with coefficients $c$ inflates the gradient's shot variance by $\|c\|_1^2 = 4$ at equal shots per circuit~\cite{Banchi2026}; worst-case bounds~\cite{Quek2024} do not apply to a shallow local circuit; inflation is predicted to grow with the cone's gate count~\cite{Quek2024,Kim2023}. Fallback: if the gain-3 expectation value is within two standard errors of zero the point is reported unmitigated and flagged; the extrapolator is never switched per point. \emph{Refuted if} level 1 shifts the landscape-to-shot variance ratio by more than its bootstrap interval at any main-grid point, or the level-2 inflation at $n = 40, 100$ and $L = 2, 8$ differs from 4 by more than its interval, or does not grow with cone size from $L = 2$ to $L = 8$.

\subsection{H5--H7: the reset dial}
\textbf{H5.} $\Var_\theta[\Cmix]$ at $n = 60$ follows the pre-drawn second-moment curve at each $p$: nearly $L$-independent between $L = 8$ and 12 and above the Deviation 33 floor $\tfrac12 p^2(c_i^2 g_i^2 + c_j^2 g_j^2)$ ($2.20\times10^{-3}$ and $1.50\times10^{-2}$ at $n = 53$, cleared by the predictions by 1.62 and 1.22) after shot and pattern floors are subtracted; $p^4/9$ is a historical reference line only. \emph{Refuted if} the paired-bootstrap interval of $\Var[\Cmix](L{=}12)/\Var[\Cmix](L{=}8)$ misses the pre-drawn ratio by more than the combined bootstrap-plus-truncation interval at either $p$, or the floor-subtracted value misses the pre-drawn value at either $L$ and $p$, or the floor-subtracted $\Var[\Cmix]$ lies below the Deviation 33 floor with its interval at either $p$; a value between the $p^4/9$ line and the ansatz floor is not a finding about anything.
\textbf{H6.} At $p \in \{0.25, 0.5\}$ the $k = L$ gradient variance follows its pre-drawn near-flat curve in $L$ and in $n$ ($n = 40, 60, 100$ at $p = 0.25$) while $k = 1$ decays with $L - k$; the unital references (delay-matched $p = 0$, dephasing dial) fall. The $k = L$ variance is bounded below, rigorously for this ansatz and channel, by $\tfrac12 c_i^2 g_i^2 p^2$ (Deviation 33; $1.06\times10^{-3}$ and $7.35\times10^{-3}$ at $n = 53$, cleared by 1.88 and 1.30), $L$- and $n$-independent; the $k = 1$ series has no such floor, and that contrast is the content of H6. Statistics (Deviation 40): the $k = L$ depth ratio $V(12)/V(8)$, predicted 1.00 for the dial against 0.023--0.044 for the unital references, and the absolute $k = L$ separation from the unital references; the dial's $k = 1$ rows lie below the shot floor by prediction and are reported as upper bounds, so no dial ratio with a $k = 1$ denominator is formed. \emph{Refuted if} the floor-subtracted $k = L$ variance at any grid or ladder point lies below the Deviation 33 floor with its interval, or the $V(12)/V(8)$ ratio at $k = L$ misses the pre-drawn ratio at either $p$, or the ladder ratio $\Var(n{=}100)/\Var(n{=}40)$ at $p = 0.25$ misses its pre-drawn value, or the reset dial's $k = L$ variance at $n = 60$, $L = 8$ does not exceed the dephasing dial's by the pre-drawn factor, or the $k = 1$ series does not fall with $L$. A flat unital reference on the day makes the ladder comparison inconclusive rather than refuting. On-day $M$ rule (Deviation 39): if the measured bootstrap $2\sigma$ of the $L = 8$, $p = 0$ reference on a rung exceeds half its predicted depth fall, $M$ rises from 350 to 600 on that rung at both depths; the estimand is unchanged.
\textbf{H7.} At $p = 0.5$, $n = 60$, $L = 8$ the RMS over draws of $\Cmix - \Cmix[L-\ell, L]$ falls with $\ell$ as $c^{\ell/2}\,\mathrm{std}_\theta(\Cmix)$ up to an order-one constant from the pre-drawn curve; $\ell = 2$ is resolved at $2.5\sigma$ at the $p^4/9$ bound (better from the Deviation 33 floor, $\mathrm{std}(\Cmix) \geq 0.122$), $\ell = 4$ is an upper-bound point unless the measured $\mathrm{std}(\Cmix) > 0.1$. \emph{Refuted if} the $\ell = 2$ RMS misses the pre-drawn prediction by more than the combined interval, or, when $\mathrm{std}(\Cmix) > 0.1$, the $\ell = 4$ RMS is not below the $\ell = 2$ RMS by more than the paired interval.

\subsection{Kill rules and anomaly protocol}
Gate 2 pause rule: failing any Gate 2 criterion pauses the campaign until the cause is logged and the PI signs a Deviation. Dial kill rule: the arm is not run if the measured reset error exceeds $2\times10^{-2}$ on any dial-patch qubit that cannot be swapped out at placement (Deviation 42 (vii)), QPU-locked time exceeds 7.0\,min per dial gradient point including job overhead at the submitted $\tau_{\mathrm{rep}}$ (the Deviation 27 budget of 5.95\,min with 18\% headroom; Deviation 41; the circuit-execution-only time is reported beside it), the transpiler maps reset to measure-plus-conditional-$X$ (a dial layer above 1\,$\mu$s, taken as the 0.35\,$\mu$s layer plus the target's reset duration; Deviation 42 (vi)), or the Estimator refuses a mid-circuit reset at the grid's resilience level. Gate 2 (c)'s ``twice the 4096-shot floor'' on a probability is $2\times1/(2\sqrt N) = 1.56\times10^{-2}$ (Deviation 42 (iv)); Gate 2 (a) uses the simulated null floor of criterion (d) until a null-control point type exists in the job-list schema (candidate Deviation 43; Deviation 42 (v)). Anomalies (Deviation 19): a single point $> 3\sigma$ from its Gate 1 prediction or a same-sign monotone deviation over three or more adjacent points in $n$ or $L$, applied literally on signed $z$-scores with the minimum $|z|$ of the run reported (Deviation 42 (viii)); confirmation requires re-measurement on another day and patch and failure of the calibrated simulations to reproduce it; unreplicated anomalies are labelled exploratory.

\section{Results}\label{sec:results}
All tables are keyed to the job-list point identifier (\texttt{patch}, $n$, \texttt{edge}, $L$, $k$, resilience, shots, $M$), the fields of \texttt{points[]} in \texttt{data/joblists/}; job tags follow \texttt{L\{resilience\}}. No value in this section exists yet.

\begin{table}[t]
\caption{Provenance. \label{tab:provenance}}
\begin{ruledtabular}
\begin{tabular}{ll}
Pre-registration version at booking & v0.11.1 (Deviations 14--42) \\
Predictions commit (\texttt{data/predictions/}) & \placeholder{sha} \\
Hardware job lists and pre-flight review links & \placeholder{list} \\
Hardware days and calibration snapshots & \placeholder{dates, files} \\
Post-run review record & \placeholder{link} \\
Analysis code commit & \placeholder{sha} \\
\end{tabular}
\end{ruledtabular}
\end{table}

\subsection{Main grid: variance versus $n$ and $L$}
\begin{table*}[t]
\caption{Main grid, $k = 1$, resilience 0 and 1. One row per job-list point. Floors: analytic two-term shot floor from the measured spreads; null-control floor from control (a). Predictions from \texttt{gate1\_predictions.csv} and \texttt{pauliprop\_predictions.csv} (Deviation 15 status shown). \label{tab:maingrid}}
\begin{ruledtabular}
\begin{tabular}{llllllllllll}
patch & $n$ & edge & $L$ & res. & $M$ & $\Var[\partial C]$ [95\% CI] & shot floor & null floor & $\eN$ & pred.\ noiseless / unital / non-unital & Dev.\ 15 \\
\hline
4x5 & 20 & 93\_103 & 1 & 0 & 200 & \placeholder{} & \placeholder{} & \placeholder{} & \placeholder{} & 2.421e-1 / 2.252e-1 / 2.260e-1 & exact \\
4x5 & 20 & 93\_103 & 2 & 0 & 200 & \placeholder{} & \placeholder{} & \placeholder{} & \placeholder{} & 8.268e-2 / 7.269e-2 / 7.246e-2 & exact \\
4x5 & 20 & 93\_103 & 4 & 0 & 200 & \placeholder{} & \placeholder{} & \placeholder{} & \placeholder{} & 1.322e-2 / \tbd{} / \tbd{} & exact / deferred \\
4x5 & 20 & 93\_103 & 8 & 0 & 200 & \placeholder{} & \placeholder{} & \placeholder{} & \placeholder{} & 3.160e-4 / 2.014e-4 / 2.156e-4 & converged \\
4x5 & 20 & 93\_103 & 12 & 0 & 200 & \placeholder{} & \placeholder{} & \placeholder{} & \placeholder{} & 1.428e-5 / 6.729e-6 / 5.240e-6 & not converged; exploratory (Dev.\ 37) \\
\multicolumn{12}{l}{\tbd{rows for $n = 39, 53, 70, 87$ at $L = 1, 2, 4, 8, 12$ and resilience 1, generated from the job lists; predictions copied from the CSVs at the provenance commit; the $n = 39$ rows are to be recomputed on edge 93\_103 (Deviation 36); every $L = 12$ row is exploratory (Deviation 37)}} \\
\end{tabular}
\end{ruledtabular}
\end{table*}

\begin{figure}[t]
\figstub{figures/variance\_vs\_n\_per\_L.pdf}{$\Var[\partial C]$ against $n$, one panel per $L$; hardware (resilience 0 and 1), noiseless, unital and non-unital predictions with truncation bands; shot floor and null-control floor on every panel; log scale.}
\caption{Gradient variance against qubit count at fixed depth. \placeholder{caption from data}}\label{fig:var_n}
\end{figure}

\begin{figure}[t]
\figstub{figures/variance\_vs\_L\_per\_n.pdf}{$\Var[\partial C]$ against $L$, one panel per $n$; exponential and power-law fits with AIC; $\Var \propto p^{2L}$ sanity line; floors.}
\caption{Gradient variance against depth at fixed qubit count. \placeholder{caption from data}}\label{fig:var_L}
\end{figure}

\subsection{Layer-index sweep}
\begin{table}[t]
\caption{Layer-index sweep at $L \in \{8, 12\}$, $n \in \{40, 100\}$ nominal, $k \in \{1, L/2, L-1, L\}$, resilience 1. $R_m$ per Deviation 14, formed only where the $k = 1$ denominator lies above the shot floor (Deviation 40); $D$ with paired-bootstrap interval; $L = 12$ rows exploratory (Deviation 37). Pre-drawn point estimates of $D$ at $L = 8$ are $|D| \leq 0.02$ with independent $2\sigma$ 0.06--0.09 (Gate 1 report, Section 6.3). \label{tab:layer}}
\begin{ruledtabular}
\begin{tabular}{lllllll}
$n$ & $L$ & $k$ & $\Var[\partial C]$ [CI] & noiseless pred. & $R_{\mathrm{hardware}}$ [CI] & $D$ [CI] \\
\hline
\multicolumn{7}{l}{\placeholder{16 rows}} \\
\end{tabular}
\end{ruledtabular}
\end{table}

\subsection{Controls and mitigation}
Null control (a): \placeholder{$\Var_{\mathrm{null}}$ per $n$ against $1/(2N)$; pre-drawn ratios 0.66--0.85}. Rotation-survival identity: \placeholder{pass count over 5 draws per point}. Mean null test: \placeholder{points with $|\bar{g}| > 2$\,s.e.}. Day-2 repeat of the $n = 40$ ladder: \placeholder{drift table}. Level 1 ratio shift: \placeholder{}. Level 2 inflation at $(n, L) \in \{40, 100\}\times\{2, 8\}$: \placeholder{measured factor against 4, against cone gate count}.

\subsection{Reset dial}
\begin{table*}[t]
\caption{Dial arm. Pre-drawn values from \texttt{pauliprop\_predictions.csv} (\texttt{stage = dial}, \texttt{gate1b}) at the $6\times10$ patch, provisional until the whole-layer $ZZ$ model (Deviation 34) is in the committed predictions. Floors: shot; pattern ($K = 256$, the independent-mask upper bound $\Var_{\mathrm{mask}}[C]/(2K)$, replaced by the per-mask-difference estimate on the data, Deviation 38); the Deviation 33 ansatz floor in the last column. The headline statistic is measured $k = L$ variance / ansatz floor (pre-drawn 1.88 at $p = 0.25$, 1.30 at $p = 0.5$). \label{tab:dial}}
\begin{ruledtabular}
\begin{tabular}{llllllllll}
arm & $p$ & $n$ & $L$ & $k$ & measured $\Var$ [CI] & floors (shot / pattern) & pre-drawn $\pm$ error & $\Var[\Cmix]$ meas.\ / pre-drawn & Dev.\ 33 floor ($k{=}L$ / $C$); meas./floor \\
\hline
dial & 0.25 & 53 & 8 & 8 & \placeholder{} & 1.22e-4 / 2.81e-4 & 1.993e-3 $\pm$ 2.6e-5 & \placeholder{} / 3.563e-3 & 1.06e-3 / 2.20e-3; \placeholder{} (pred.\ 1.88) \\
dial & 0.25 & 53 & 12 & 12 & \placeholder{} & 1.22e-4 / 2.81e-4 & 1.989e-3 $\pm$ 2.6e-5 & \placeholder{} / 3.557e-3 & 1.06e-3 / 2.20e-3; \placeholder{} \\
dial & 0.5 & 53 & 8 & 8 & \placeholder{} & 1.22e-4 / 6.71e-4 & 9.515e-3 $\pm$ 3.6e-5 & \placeholder{} / 1.821e-2 & 7.35e-3 / 1.50e-2; \placeholder{} (pred.\ 1.30) \\
dial & 0.5 & 53 & 12 & 12 & \placeholder{} & 1.22e-4 / 6.71e-4 & 9.515e-3 $\pm$ 3.6e-5 & \placeholder{} / 1.821e-2 & 7.35e-3 / 1.50e-2; \placeholder{} \\
delay-matched & 0 & 53 & 8 & 8 & \placeholder{} & 3.05e-5 (16384 shots, Dev.\ 35) / -- & 2.821e-4 $\pm$ 7.7e-6 & \placeholder{} / 2.843e-4 & -- \\
delay-matched & 0 & 53 & 12 & 12 & \placeholder{} & 1.22e-4 / -- & 6.404e-6 $\pm$ 1.0e-6 & \placeholder{} / 6.822e-6 & -- \\
dephasing & 0.5 & 53 & 8 & 8 & \placeholder{} & 1.22e-4 / -- & 1.705e-5 $\pm$ 2.1e-6 & \placeholder{} / 1.836e-5 & -- \\
ladder & 0.25 & 39 & 8 & 8 & \placeholder{} & 1.22e-4 / 2.88e-4 & 2.094e-3 $\pm$ 2.7e-5 (edge 94\_95; to be recomputed on 93\_103, Dev.\ 36) & \placeholder{} / \tbd{} & \tbd{per patch} \\
ladder & 0.25 & 87 & 8 & 8 & \placeholder{} & 1.22e-4 / 2.83e-4 & 2.028e-3 $\pm$ 2.6e-5 & \placeholder{} / \tbd{} & \tbd{per patch} \\
\multicolumn{10}{l}{\tbd{$k = 1$ points (upper bounds, Deviation 40; contingent), $p = 0$ ladder references at $M = 350$, or 600 where the on-day rule of Deviation 39 fires (n = 39/53/87, $L = 8$ and 12; $n = 53$, $L = 8$ at 16384 shots), truncation arm $\ell = 2, 4$}} \\
\end{tabular}
\end{ruledtabular}
\end{table*}

\begin{figure}[t]
\figstub{figures/dial\_variance\_vs\_p.pdf}{$k = L$ gradient variance and $\Var[\Cmix]$ against $p \in \{0, 0.25, 0.5\}$ at $n = 60$, $L = 8$ and 12, with the delay-matched and dephasing controls, the pre-drawn second-moment curves with truncation bands, the $p^4/9$ reference line, and the shot and pattern floors.}
\caption{Reset dial. \placeholder{caption from data}}\label{fig:dial}
\end{figure}

\begin{figure}[t]
\figstub{figures/prediction\_vs\_measured.pdf}{Measured against pre-drawn variance for every point of the grid, the sweep and the dial, coloured by model; identity line; Deviation 15 hardware-only points open symbols; anomalies (Deviation 19) marked.}
\caption{Prediction against measurement. \placeholder{caption from data}}\label{fig:pred_meas}
\end{figure}

Truncation arm (H7): \placeholder{RMS at $\ell = 2, 4$ after residual-pattern subtraction, against $c^{\ell/2}\,\mathrm{std}(\Cmix)$}. Fitted exponent $\alpha(p)$ at $p \in \{0.25, 0.5\}$ against $\log(1/c(p))$: \placeholder{}. Headline: measured $k = L$ variance over the Deviation 33 floor, \placeholder{value [CI] at each $p$} against the pre-drawn 1.88 and 1.30. Free pipeline check $\mathbb{E}_\theta[\Cmix] = p^2$ folded through readout (0.066 and 0.252 predicted at the $6\times10$ patch): \placeholder{}.

\subsection{Anomalies and exploratory points}
\placeholder{List under Deviation 19: point, deviation in $\sigma$, replication attempt, simulation reproduction; or ``none''.} Points designated hardware-only (Deviation 15), exploratory ($L = 12$, Deviation 37) or inconclusive (Gate 1b branch) are labelled in every table.

\section{Discussion}\label{sec:discussion}
The pre-registration permits pre-written interpretation for each outcome branch; the paragraphs below were drafted before data and only the branch that occurred is kept at submission, the others moved to the Supplement with a note.

\paragraph{H1 and H2 pass.} \tbd{Branch text.} The depth curve is set by unital gate noise at the snapshot error rates, as the calibrated simulation predicted; the $n$ dependence at $L \leq 4$ is geometric (cone versus patch) and at $L \geq 8$ the truncated cone and broken couplers, not concentration in $n$, dominate. The measurement then adds to the literature a hardware variance curve with intervals against a curve pre-drawn from calibration, and nothing about trainability beyond what the simulation already implied.

\paragraph{H2 fails (variance does not fall with $L$ above the floor).} \tbd{Branch text.} Either the signal lies below the measured null-control floor plus 3 bootstrap $\sigma$ (criterion (d) as amended by Deviation 37) and the row is an upper bound, or the noise model misses a channel; the Deviation 19 protocol applies before any claim, and the result is reported as a discrepancy between hardware and the calibrated model, not as evidence against the theorems, which concern ensembles the ansatz does not realise.

\paragraph{H3 consistency.} $D = 0$ within the interval is the prediction; a paired interval excluding zero at native noise would be an anomaly under Deviation 19 and would need replication before interpretation.

\paragraph{H5--H6 pass.} \tbd{Branch text.} A controlled non-unital channel raises the cost-variance floor and the late-layer gradient variance while the early-layer gradient collapses, and the measured excess over the rigorous ansatz floor matches the first-principles noise model at the \placeholder{level} level: the arm is a variance-level benchmark of that model (Deviation 33). The channel is inside the class for which Angrisani et al.~\cite{Angrisani2026} prove average-case Pauli-propagation simulability, but that guarantee is over 2-design gate ensembles and, at these noise levels and gradient-scale precision, its runtime bound is not informative for our ensemble, so classical simulability of these circuits is demonstrated numerically by the baseline beside every point, not inferred from the theorem, and is not the paper's conclusion. A residual between hardware and the truncated classical baseline is a noise-model or truncation discrepancy, not evidence of hardness.

\paragraph{H5--H6 fail or inconclusive.} \tbd{Branch text.} If the delay-matched reference is flat on the day, the ladder comparison is inconclusive and is reported as such; if the $k = L$ series falls where the pre-drawn curve is flat, the candidates are an unmodelled channel (idle-branch $T_1$; a $ZZ$ residual beyond the whole-layer static model of Deviation 34), a reset that does not realise $\Np$ (Paper 2, Q4), or a pattern-noise estimate in error; the order in which these are checked is fixed here.

\paragraph{H7.} \tbd{Branch text for resolved / upper-bound $\ell = 4$.}

\paragraph{H4.} \tbd{Branch text: level-1 ratio unchanged as algebra predicts; level-2 inflation against $\|c\|_1^2 = 4$ and cone size; this is the first average-case hardware test of a worst-case bound, framed as such~\cite{Larocca2025,Quek2024}.}

\paragraph{Limits.} One differentiated parameter per point; loss-first diagnosis is served by $\Var[C]$~\cite{Larocca2025,Arrasmith2022}. Two values of $p$ do not test monotonicity of $\alpha(p)$ in $\log(1/c)$. Ry/CZ is not a local 2-design, so no constant of the theorems is tested, only functional forms against ansatz-specific second-moment curves. Gate 1b clause (b) was amended three times before data and clarified once more (Deviation 35). The $L = 12$ rung is below the 4096-shot floor in every prediction and is exploratory (Deviation 37); the per-row 4096-versus-16384 resolvability table at $L = 8$ was \tbd{computed before Gate 2; candidate Deviation 44 if any row moved}.

\section*{Data availability}
Repository: \url{https://github.com/SSIT-Q/gradvar-phoenix} (Apache 2.0): code, job lists with pre-flight review links, one metadata bundle per job (job metadata, usage, full result, options, backend calibration and target snapshot, transpiled circuits), the calibration snapshots in \texttt{data/calibrations/}, the predictions in \texttt{data/predictions/}, the pre-registration with its Deviations table, and the review notes. Zenodo DOI: \tbd{minted at submission}. The IBM API key is read only from repository secrets and never appears in the repository.

\section*{Author contributions}
\tbd{Name}: design, code, runs, analysis, first draft. \tbd{Name}: PI, supervision, arXiv endorsement. \tbd{Name or acknowledgement}: theory comparison, review of the analysis plan. Independent code and pre-flight reviews: \tbd{as recorded in the repository}.

\section*{Acknowledgements}
\tbd{Internal referee (anonymity honoured); IBM Quantum Flex plan; college.}

\appendix

\section{Deviations from the signed pre-registration}\label{app:deviations}
Table~\ref{tab:deviations} is generated from \texttt{docs/manuscripts/deviations.yaml} by \texttt{scripts/build\_deviations\_table.py}; rows 1--13 of the pre-registration (citation corrections between v0.1 and v0.2) are summarised in the pre-registration itself. \tbd{Any Deviation recorded after v0.11.1 (candidates 43, null-control point type; 44, $L = 8$ shot-resolvability table) is appended to the YAML, never to this file.}
\onecolumngrid
% Auto-generated by scripts/build_deviations_table.py from deviations.yaml on 2026-09-20. Do not edit by hand.
\begin{longtable}{@{}p{0.05\linewidth}p{0.11\linewidth}p{0.20\linewidth}p{0.42\linewidth}p{0.16\linewidth}@{}}
\caption{Deviations from the signed Paper 1 pre-registration, all recorded before any Paper 1 datum (rows 14--42; rows 1--13 are the citation corrections of v0.2; row 32 was recorded before the smoke test's data existed and fixes how they may be used). Source: preregistration\_q1.html v0.11.1 (19-20 Sep 2026), Section 9.}\label{tab:deviations}\\
\toprule
No. & Date & Title & Content & Approval \\
\midrule
\endfirsthead
\toprule
No. & Date & Title & Content & Approval \\
\midrule
\endhead
\bottomrule
\endfoot
14 & 2026-09-19 & H3 test statistic & The dimensionally inconsistent ratio-versus-variance comparison is replaced by the noiseless-corrected layer-index ratio $R_m = [\mathrm{Var}_m(k{=}L)/\mathrm{Var}_m(k{=}1)] / [\text{same, noiseless}]$ on shared draws; H3 supported if the paired-bootstrap 95\% interval of $R_{\mathrm{hardware}}/R_{\mathrm{unital}}$ excludes 1 in the Mele direction. & Approved: Dr. Raviram V, 19 Sep 2026 \\
15 & 2026-09-19 & Gate 1 predictions at n = 40 and 100 for L = 8 and 12 & Exact noisy simulation infeasible at those points; predictions come from truncated Pauli propagation with a reported truncation error; a point whose error exceeds half the unital-vs-noiseless separation is hardware-only and its result exploratory. & Approved: Dr. Raviram V, 19 Sep 2026 \\
16 & 2026-09-19 & H3 reclassified & At native damping ($\|t\| \approx 2\times10^{-3}$ per layer) the predicted $|D| \lesssim 10^{-2}$ is unresolvable against \textasciitilde{}0.3 measurement uncertainty at M = 200; H3 becomes a consistency check with prediction D = 0; the layer-index sweep is retained as a check of the unital k-dependence (H2). & Approved: Dr. Raviram V, 19 Sep 2026 \\
17 & 2026-09-19 & Gate 1 criterion (b) bound & Bootstrap hi/lo bound made depth-dependent for heavy-tailed gradients (kurtosis 2.3 / 4.2 / 8.4 at $L = 1 / 2 / 4$): $< 1.5$ at $L = 1$, $< 2.0$ at $L = 2$, $< 2.5$ at $L \geq 4$ with $M = 200$; M raised to 400 (L = 2) and 700 (L >= 4) where shots allow, with the original 1.5 bound. & Approved: Dr. Raviram V, 19 Sep 2026 \\
18 & 2026-09-19 & Ladder qubit counts & Readout cut excludes qubit 107, so the 4x10 point is n = 39; ladder recorded as 20 / 39 / 56 / 71 / 90 on the 19 Sep snapshot, re-derived on the run day; recorded 20 Sep under Deviations 22 and 26 as 20 / 39 / 53 / 70 / 87 (6x10 at origin (6,0), 8x10 at (4,0), 10x10 at (2,0); edges 84\_85, 84\_85, 75\_85). & Approved 19 Sep; recount countersigned Dr. Raviram V, 20 Sep 2026 \\
19 & 2026-09-19 & Anomaly protocol & An anomaly is a $> 3\sigma$ single-point deviation from its Gate 1 prediction or a same-sign monotone deviation over three adjacent points; confirmation needs re-measurement on another day and patch ($\leq 20$ reserve minutes) and failure of the calibrated simulations to reproduce it; unreplicated anomalies are reported as exploratory. & Approved: Dr. Raviram V, 19 Sep 2026 \\
20 & 2026-09-19 & Non-unital arm added (hybrid plan) & Section 3b added: a controlled reset dial realising $N_p(\rho) = (1-p)\rho + p\,|0\rangle\langle 0|$ through the native 400 ns reset with fresh random masks per layer and circuit, estimator $C_{\mathrm{mix}}$, hypotheses H5-H7 as second-moment predictions for this ansatz, Gate 1b, matched controls, truncation arm, kill criteria; simulability stated as the conclusion. & Approved: Dr. Raviram V, 19 Sep 2026 \\
21 & 2026-09-19 & Corollary 6 floor corrected to $p^4/9$ & The Mele et al. cost-variance lower bound is $(\|t\|_2^2/3)^{|P|} = p^4/9$ for $|P| = 2$ (4.3e-4 at p = 0.25, 6.9e-3 at p = 0.5), not $(1/3)\|t\|_2^{2|P|}$; H5 refutes against the pre-drawn ansatz curve with the 2-design bound as a reference line; H7 bounds and the Section 3b review items updated. & Adopted under delegated authority; countersigned Dr. Raviram V, 20 Sep 2026 \\
22 & 2026-09-20 & Exclusion list extended from raw properties & Exclusion rule gains $|ZZ| \geq 1$ MHz to an excluded or dead qubit and initialisation error $\geq 5\times10^{-4}$, both from the raw backend properties; on the 19 Sep snapshot this excludes 18, 27 (ZZ) and 8, 11, 22, 59 (init error); per-edge ZZ values enter the idle-phase simulation model. & Adopted under delegated authority; countersigned Dr. Raviram V, 20 Sep 2026 \\
23 & 2026-09-20 & Repetition-delay feasibility gate & Gate 2 gains a rep\_delay criterion: default\_rep\_delay, rep\_delay\_range and dynamic\_reprate\_enabled read from backend.configuration(); grid booked at a named floor (1 us or the smallest bias-free value on the smoke-test ladder); pre-registered cut order if the grid exceeds 200 min; ibm\_phoenix default measured at 1.0 us on 20 Sep 2026. & Adopted under delegated authority; countersigned Dr. Raviram V, 20 Sep 2026 \\
24 & 2026-09-20 & Budget model corrected from the Marrakesh pipeline check & Per-job model $T = 2\,\mathrm{s} + (\tau_{\mathrm{rep}} + L\times0.71\,\mu\mathrm{s} + t_{\mathrm{meas}} + 10\,\mu\mathrm{s}\,[+\text{dial}])\,N_{\mathrm{exec}}\,(3\text{ at resilience 2}) + [\text{resilience}\geq1]\,32\,\text{shots}\,n_{\mathrm{bases}}\,(\tau_{\mathrm{rep}} + t_{\mathrm{meas}} + 10\,\mu\mathrm{s})$; back-predicts 21.1 s against 22 s measured; tables recomputed (Paper 1 grid \textasciitilde{}90 min at 1 us; dial core 57.8 min, full 85.3 min); runner layout re-check added. & Adopted under delegated authority; countersigned Dr. Raviram V, 20 Sep 2026 \\
25 & 2026-09-20 & Criterion (a) estimator amendment & The chain regression check cannot be met at $n \geq 13$ by a percentile bootstrap (kurtosis $(3/2)^n$); amended to (a-i) per-draw closed-form identity $< 10^{-10}$ at $n = 4..20$, (a-ii) bootstrap coverage at $n \leq 12$ ($M = 1000$), (a-iii) exact null sampling distribution at $n = 13..20$ ($M = 300$), with a seed-2027 replicate rule. & Adopted under delegated authority; countersigned Dr. Raviram V, 20 Sep 2026 \\
26 & 2026-09-20 & Pre-submission layout re-check & Before every submission the runner checks the placed patch against the day's properties: readout $< 3\times10^{-2}$, CZ error $< 5\times10^{-3}$ per coupler (couplers above the cut are broken edges carrying no CZ; the observable edge must be intact), init error $< 5\times10^{-4}$, the Deviation 22 ZZ rule; verdict logged (layout\_check); override never used on the observable edge or its L = 2 cone. & Adopted under delegated authority; countersigned Dr. Raviram V, 20 Sep 2026 \\
27 & 2026-09-20 & Mixture-estimator mask count & K = 256 masks x 16 shots per (draw, shift) replaces K = 64 x 64: the pattern-noise floor $\mathrm{Var}_{\mathrm{mask}}[C]/(2K)$ falls from $\approx1.1\times10^{-3}$ to $\approx2.8\times10^{-4}$ so Gate 1b(b) and (d) pass; 51,200 circuits and 171 jobs per gradient point; dial line re-booked from 45 to 60 min. & Adopted under delegated authority; countersigned Dr. Raviram V, 20 Sep 2026 \\
28 & 2026-09-20 & Gate 1b clause (b): per-series floors and M = 500 for the p = 0 ladder & The delay-matched p = 0 points carry no mask lottery, so their floor is the shot floor alone; clause (b) amended to a 3x-shot-floor and 2x-draw-2sigma test on the n-ladder fall; p = 0 ladder points booked at M = 500; dial line 65 min. Superseded in part by Deviation 29 the same day. & Adopted under delegated authority; countersigned Dr. Raviram V, 20 Sep 2026 \\
29 & 2026-09-20 & Gate 1b clause (b): unital reference resolvability as a depth fall & With the measured kurtosis 14.2 the n-ladder fall is unresolvable at any affordable M and non-monotone in n (broken couplers); clause (b) becomes the L = 8 -> 12 fall of the delay-matched p = 0 k = L variance on each patch, $> 3\times$ shot floor and $> 2\times$ the $L = 8$ draw $2\sigma$; three p = 0 L = 12 points added; M = 500 reverts to 200. & Adopted under delegated authority; specifically countersigned Dr. Raviram V, 20 Sep 2026 \\
30 & 2026-09-20 & Gate 1b clause (b): draw count from the measured kurtosis; per-rung evaluation & Clause (b) evaluated per rung, passing if $\geq 2$ of 3 rungs pass, a rung below 3 shot floors at L = 8 being 'reference unresolvable'; p = 0 reference points booked at M = 350 by the rule $M = 17.5(1.3\kappa - 1)$ (second-review correction); ZZ idle phase to be added to the propagation model before Gate 2; clause (b) frozen. & Adopted under delegated authority; specifically countersigned Dr. Raviram V, 20 Sep 2026 \\
31 & 2026-09-20 & Schedule brought forward & Hardware window moved from 11-25 Oct to 27 Sep-12 Oct 2026: Gate 1 decision 23-26 Sep, ibm\_phoenix smoke test 27-29 Sep (\textasciitilde{}5 Flex min), Gate 2 the same week, main grid, dial arm and Paper 2 characterisation 1-12 Oct; no hypothesis, threshold, grid point, shot count or budget line changes; gates decided in order. & Adopted under standing delegated authority; countersignature on the PI's next review \\
32 & 2026-09-20 & Smoke test brought forward to 20 Sep 2026 & Job list 02 (six grid points at $n = 20$, $L = 2$ and 8 at resilience 0, 1, 2; dial and delay-matched probes; reset-error sequence; rep\_delay ladder 1 / 5 / 20 / 250 $\mu$s) runs on 20 Sep 2026, about 0.7 Flex min at 1 $\mu$s from the dry-run line; its data feed Gate 2 criteria (a)-(e) and the Section 3b kill rules only, never Gate 1 or any hypothesis; 27-29 Sep held for a re-run. & Adopted under standing delegated authority at the co-author's request; countersignature on the PI's next review \\
33 & 2026-09-20 & Ansatz-specific variance floor & The 2-design bound $p^4/9$ is replaced as the H5/H6 refutation threshold by floors proven for this ansatz and channel from uniform angles alone: $\mathrm{Var}[\partial_{k=L}C] \geq \tfrac12 c_i^2 g_i^2 p^2$ and $\mathrm{Var}[C] \geq \tfrac12 p^2(c_i^2 g_i^2 + c_j^2 g_j^2)$, $L$- and $n$-independent, pre-drawn per patch and estimator (at $n = 53$: 1.06e-3 / 7.35e-3 for $k = L$, 2.20e-3 / 1.50e-2 for $C$ at $p = 0.25$ / 0.5; predictions clear them by 1.88 / 1.30 and 1.62 / 1.22); measured $k = L$ variance over this floor is the Section 3b headline; the arm is a variance-level benchmark of a first-principles noise model, not a test of Theorem 8 nor a demonstration of simulability; $p^4/9$ kept as a historical line. & Adopted under delegated authority (theorist role delegated 20 Sep 08:10 IST); theorist and PI countersignature on return \\
34 & 2026-09-20 & Whole-layer static ZZ in both propagation models & Both the main-grid unital model and the dial model carry, for every lattice pair with a reported coupling, $e^{-i\zeta\tau Z_aZ_b/4}$ per layer over the scheduled layer time outside calibrated CZs (up to 0.35 $\mu$s Paper 1, 0.71 $\mu$s dial layer), propagated by the incoherent rule moving weight $\sin^2(\zeta\tau/2)$ per pair per layer (8.8e-4 to 3.6e-3 at the 27 kHz median); convention $\zeta = 2\pi\times$ the reported coupling, conditional phase $\zeta\tau$; validated on a 2x4 density-matrix check; readout folding assumes uncorrelated confusion, checked against Paper 2 Q1 (d). & Adopted under delegated authority (theorist role delegated 20 Sep 08:10 IST); theorist and PI countersignature on return \\
35 & 2026-09-20 & Gate 1b clause (b) clarification, pre-data & A rung's $L = 8$, $p = 0$ reference is judged resolvable on the measured, floor-subtracted value at its shot count, never on the prediction; clause (b) passes only if at least two rungs are counted and all counted rungs pass, and is inconclusive if fewer than two are counted; the $n = 53$, $L = 8$, $p = 0$ reference runs at 16384 shots (floor 3.05e-5; predicted 9.2 floors, fall 9.0 floors) so all three rungs can count, +2.7 min at 1 $\mu$s from the reserve; no change of direction or power. & Adopted under delegated authority; marked for the PI's specific countersignature on return \\
36 & 2026-09-20 & 4x10 observable edge and cone-graph statement & The 4x10 observable edge moves from (94, 95) to (93, 103) because the broken coupler (95, 96) sat on an observable qubit inside every cone from $L = 2$; the light-cone statement becomes cone-graph independence (exactly $n$-independent at $L = 1$; at $L \geq 2$ only between rungs with isomorphic cone graphs after holes and broken couplers: 4x5 and 4x10 at $L = 2$, 6x10 and 8x10 at $L = 2$; the 10x10 edge (75, 85) matches no rung; the earlier claim that 87 matches at $L \leq 4$ is withdrawn); a same-geometry curve tests the noise model on that graph, not $n$-dependence; the $n = 39$ Gate 1 and Gate 1b rows are recomputed on the new edge before the Gate 1 decision. & Adopted under delegated authority (theorist role delegated 20 Sep 08:10 IST); theorist and PI countersignature on return \\
37 & 2026-09-20 & L = 12 rows exploratory; criterion (d) allowance & At $M = 350$ the variance-estimate $2\sigma$ (1.8e-5) exceeds every $L = 12$ unital prediction, and 32768-65536 shots would be needed: the main-grid $L = 12$ rows and the dial $k = 1$ rows at $L = 12$ are exploratory and enter no confirmatory statistic (not-converged propagation rows labelled and in no ratio); dial $k = L$ and $p = 0$ references at $L = 12$ keep their roles; Gate 1 criterion (d)'s fixed 10x allowance is replaced by: the smallest signal to be claimed lies above the measured null-control floor plus 3 bootstrap $\sigma$; $L = 8$ keeps 4096 shots, rows below the bar reported as upper bounds; the 4096-vs-16384 resolvability table is computed before Gate 2 (candidate Deviation 44). & Adopted under delegated authority (theorist role delegated 20 Sep 08:10 IST); theorist and PI countersignature on return \\
38 & 2026-09-20 & Dial masks: shared or independent between shift circuits & The two shift circuits of a draw may share their mask sequence or draw independent ones, fixed per point before submission and logged (mask\_seed); with shared masks the gradient pattern floor falls below $\mathrm{Var}_{\mathrm{mask}}[C]/(2K)$, which becomes the independent-mask upper bound; in both cases the gradient pattern floor is estimated directly from the $K$ per-mask gradient differences (variance over masks minus shot variance, divided by $K$); $\mathrm{Var}_\theta[C_{\mathrm{mix}}] = \mathrm{Var}_\theta[C_{\Phi_{\mathrm{mix}}}] + \mathbb{E}_\theta\mathrm{Var}_{\mathrm{mask}}[C]/K$ exactly. & Adopted under delegated authority (theorist role delegated 20 Sep 08:10 IST); theorist and PI countersignature on return \\
39 & 2026-09-20 & On-day M rule for the p = 0 references & If the measured bootstrap $2\sigma$ of the $L = 8$, $p = 0$ reference on a rung exceeds half its predicted depth fall, $M$ rises from 350 to 600 on that rung at both depths (about 1.5 min per rung at 1 $\mu$s, up to 4.5 min itemised in the reserve); estimand unchanged, no log-variance or trimmed estimator substituted; a trigger on the sample kurtosis itself is rejected (20-30 percent error). & Adopted under delegated authority (theorist role delegated 20 Sep 08:10 IST); theorist and PI countersignature on return \\
40 & 2026-09-20 & Deviation 14 ratio scope; H6 statistic & The Deviation 14 ratio $R_{\mathrm{hardware}}/R_{\mathrm{unital}}$ is reported only where the $k = 1$ denominator lies above the shot floor (predicted 0.98-1.05 with $2\sigma$ 0.04-0.42 at native noise: a consistency check with no power); under the dial the $k = 1$ variance is below the shot floor by prediction, so those rows are upper bounds and no dial ratio with a $k = 1$ denominator is formed; the H6 statistic is the $k = L$ depth ratio $V(12)/V(8)$ (predicted 1.00 for the dial against 0.023-0.044 for the unital references) together with the absolute $k = L$ separation from the unital references. & Adopted under delegated authority (theorist role delegated 20 Sep 08:10 IST); theorist and PI countersignature on return \\
41 & 2026-09-20 & Kill rule 3b (b) made consistent with Deviation 27 & The kill line becomes 7.0 min of QPU-locked time per dial gradient point including job overhead at the submitted rep\_delay (the Deviation 27 budget of 5.95 min with 18 percent headroom, matching the 65-min dial line), replacing the 5-min line that the booked design would have exceeded at any execution speed; the circuit-execution-only time (15.3 s per point at 1 $\mu$s) is reported beside it and both are extrapolated from the smoke test's per-execution and per-job constants; rules (a), (c), (d) unchanged. & Adopted under delegated authority (theorist role delegated 20 Sep 08:10 IST); theorist and PI countersignature on return \\
42 & 2026-09-20 & Pre-data analysis clarifications & Ten readings the analysis code implements: (i) H1 part 2 refers to the measured $L = 12$ curve, reported not decided; (ii) H2 refuted by any evaluable series or $k = L$ pair, all listed; (iii) H3 refutation is the Deviation 14 clause, Deviation 16 reading alongside, Deviation 40 scope; (iv) Gate 2 (c) probability floor $2\times1/(2\sqrt N) = 1.56\times10^{-2}$; (v) Gate 2 (a) uses the simulated null floor until a null-control point type exists (candidate Deviation 43); (vi) kill rule 3b (c) dial-layer duration = 0.35 $\mu$s layer + the target's reset duration; (vii) kill rule 3b (a) lists every dial-patch qubit above 2e-2, swappability decided at placement; (viii) Deviation 19 (ii) applied literally on signed $z$-scores of adjacent points, minimum $|z|$ reported; (ix) Deviation 25 on hardware: approximate null by resampling, intervals labelled approximate; (x) H4 and the fits exclude unclaimable points ($\varepsilon_N \geq 1$), listed. Renumbering: the $L = 8$ shot-resolvability table becomes candidate Deviation 44. & Adopted under delegated authority (theorist role delegated 20 Sep 08:10 IST); theorist and PI countersignature on return \\
\end{longtable}

\twocolumngrid

\section{Second-moment Pauli propagation}\label{app:pauliprop}
Summary of \texttt{docs/PAULIPROP.md}; the full derivation is the theory companion's first task. Back-propagating $O = Z_iZ_j$ (readout-folded) through the transpiled circuit ($R_y(\theta) = R_z\,\sqrt{X}\,R_z(\pi+\theta)\,\sqrt{X}\,R_z$ in the $\{R_z, \sqrt{X}, X, \mathrm{CZ}\}$ basis, calibration error after every $\sqrt{X}$ and CZ) writes $\langle O\rangle(\theta) = \sum_{\mathrm{paths}} c_{\mathrm{path}}(\theta)$ with $c_{\mathrm{path}} = \mathrm{const}\times\prod_g f_g(\theta_g)$, $f_g \in \{1, \cos, \sin\}$ per rotation. For uniform angles $\mathbb{E}[\cos] = \mathbb{E}[\sin] = \mathbb{E}[\cos\sin] = 0$ and $\mathbb{E}[\cos^2] = \mathbb{E}[\sin^2] = 1/2$; two distinct paths differ in factor type at some gate and have zero covariance, so
\[
\mathbb{E}[\langle O\rangle^2] = \sum_{\mathrm{paths}} 2^{-\#\mathrm{branchings}}\,(\text{noise factors})^2\,[P_{\mathrm{final}} \in \{I, Z\}^{\otimes n}],
\]
and the derivative keeps only paths whose Pauli fails to commute with the generator at gate $(k, q)$, with weight $1/2$; $\mathbb{E}[\partial\langle O\rangle] = 0$. The second moment is therefore a positive linear map on weights over Pauli strings: rotation about axis $A$ sends $P \notin \{I, A\}$ to $1/2$ on each of the two other non-identity Paulis; Cliffords permute; a single-qubit channel $r \mapsto Dr + t$ with diagonal $D$ sends $P_a \mapsto D_a^2 P_a$ and $t_a^2$ to $I$; a Pauli channel multiplies non-identity strings by $(1-\lambda)^2$. The only $\theta$-independent segment is the prefix before the first rotation (the last layer's CZ-block noise and, for the dial, the last $\Np$), propagated at the coefficient level with signs and squared afterwards. Channel rules: unital = depolarising factors $(1 - 2\epsilon)$ per $\sqrt X$ and $(1 - 4\epsilon/3)$ per CZ; non-unital = thermal relaxation $D = (e^{-t/T_2}, e^{-t/T_2}, e^{-t/T_1})$, $t_z = 1 - e^{-t/T_1}$; reset dial $D = (1-p)(e^{-400/T_2}, e^{-400/T_2}, 1)$, $t_z = p$; dephasing dial $t = 0$. Two engines: deterministic truncation by coefficient (rigorous lower bound; the discarded mass is a diagnostic, not an error bound) and unbiased Pauli-path Monte Carlo ($N = 2\times10^6$ paths). Pattern-noise floor: $\mathbb{E}_\theta\Var_{\mathrm{mask}}[C] = \mathbb{E}_{\mathrm{mask},\theta}[C_{\mathrm{mask}}^2] - \mathbb{E}_\theta[\Cmix^2]$ by a sampled propagation with a fresh reset mask per path and layer. Static $ZZ$ over the whole scheduled layer outside calibrated CZs is carried in both models by the incoherent rule (weight $\sin^2(\zeta\tau/2)$ per pair per layer, rotation $e^{-i\zeta\tau Z_aZ_b/4}$, $O(p\sin^2(\zeta\tau/2))$ cross-term residual stated, validated on a $2\times4$ density-matrix check; Deviation 34). Known gaps: two $Z \to I$ relaxation branches separated by a CZ are treated as distinct paths (relative $< 10^{-3}$, positive); idle-branch $T_1$ of the dial ($t_z \approx 2\times10^{-3}$ per layer) is unmodelled and recorded in \texttt{docs/PAULIPROP.md}.

\section{Job-list identifiers and logging schema}\label{app:schema}
One row per executed circuit pair with the fields of pre-registration Section 7: \texttt{backend, job\_id, timestamp, calib\_file, n, patch\_qubits, obs\_edge, L, k, arm, p, K, mask\_seed, reset\_dur\_ns, mask\_means, mask\_se, resilience, rep\_delay, dynamic\_reprate\_enabled, layout\_check, shots, seed, theta\_hash, param\_expr, exp\_plus, exp\_minus, exp\_zero, se\_plus, se\_minus, se\_zero, zne\_gains, noise\_model\_id, gradient, depth\_t, cz\_count, fractional, notes}.

\bibliography{references}

\end{document}
